\documentclass[12pt]{article}
\usepackage[margin=0.3in,top=0.25in,bottom=0.4in]{geometry}
\makeatletter
\renewcommand{\@maketitle}{\newpage\null\vskip -3.5em\begin{center}{\LARGE \@title \par}\end{center}\par\vskip 0.2em}
\makeatother
\usepackage{amsmath,amssymb}
\usepackage{booktabs}
\usepackage{hyperref}
\setlength{\parskip}{0.6em}
\setlength{\parindent}{0pt}

\title{Why the exclusive cross section falls as $p_T^{-8}$}
\author{}
\date{}

\begin{document}
\maketitle

\section*{In short}

\begin{itemize}
\item The exclusive process is a colour-singlet exchange (two gluons). The amplitude is the
\emph{difference} of the photon wavefunction with and without the kick of the target. For a dilute
dipole, $N\propto r^2$, this difference is the \textbf{second derivative} of the wavefunction:
amplitude $\sim Q_s^2/p^4$ instead of $1/p^2$. Squared: $Q_s^4/p^8$.
\item The exact result for $N=Q_s^2r^2/4$ is Eq.~\eqref{eq:g}. The $p^{-8}$ comes from the $I_0$ term.
The $I_1$ term is $p^{-10}$.
\item $p^{-8}$ is an \textbf{asymptotic} power. In our range of $p_T$ the hard part alone is flatter
(about $p^{-6}$ to $p^{-7.4}$), because $m_c=1.5$~GeV is not small and because the dipole is not
exactly $r^2$. The photon flux, the $x_P$ dependence and the fragmentation make the full result
steeper again. The output files give a local power between $-7.5$ and $-9.7$ at $y=0$ (Section 8).
\item If the dipole is not exactly $r^2$ (BK evolution, or a logarithm), the $I_1$ term gets a
piece that falls more slowly, $p^{-2-4\gamma}$, and wins at very large $p$ (Section 6).
\end{itemize}

\section*{1. The formula in the code}

The exclusive cross section is built from two radial integrals (the function
\texttt{exclusive\_radial\_integrals} in \texttt{src/integrands.cpp}),
\begin{equation}
  I_\nu(p) = \int_0^\infty dr\; r\, J_\nu(p r)\, K_\nu(m r)\, N(r,x_P), \qquad \nu = 0,1,
\end{equation}
and the integrand is
\begin{equation}
  \frac{d\sigma}{d^2p\,dy} \;\propto\; z\, m^2\,\Big[(z^2+\bar z^2)\, I_1^2 + I_0^2\Big]
  \times \text{flux} \times \frac{D(z_h)}{z_h^2},
  \label{eq:code}
\end{equation}
with $p=p_c=p_{D^0}/z_h$ the charm transverse momentum, $z=p^+/q^+$, $\bar z=1-z$ and $m=m_c$.

Where this comes from: for a real photon ($Q^2=0$) the wavefunction of the $q\bar q$ pair has two
parts, $m\,K_1(mr)$ (the quark and antiquark have opposite helicities, weight $z^2+\bar z^2$) and
$m\,K_0(mr)$ (same helicities, needs the quark mass). The Fourier transform to the momentum $p$ of the
quark, after the angle integral, gives $J_1$ for the first and $J_0$ for the second. So $I_\nu$ is the
\textbf{amplitude} to find the quark with momentum $p$ after the dipole scattered elastically on the
target. The cross section is the amplitude squared, this is why $I_\nu$ appears squared (or as
$r_1$, $r_2$ in the 5D version).

\section*{2. At large $p$ only small dipoles matter}

$J_\nu(pr)$ oscillates with period $\sim 1/p$. Everything that is smooth in $r$ cancels, so the
integral is controlled by $r\sim 1/p$. For $p\gg Q_s$ these dipoles are small and the target is
dilute for them:
\begin{equation}
  N(r) \simeq \frac{Q_s^2 r^2}{4}.
  \label{eq:dilute}
\end{equation}
(This is the GBW form at small $r$. MV has an extra logarithm, see Section 6.)

\section*{3. The integrals without the dipole}

We need the standard integral
\begin{equation}
  \int_0^\infty dr\, r\, J_\nu(pr)\,K_\nu(mr) = \Big(\frac{p}{m}\Big)^{\nu}\frac{1}{p^2+m^2},
\end{equation}
so
\begin{equation}
  \psi_0(p)\equiv\int dr\, r\, J_0 K_0 = \frac{1}{p^2+m^2},\qquad
  \psi_1(p)\equiv\int dr\, r\, J_1 K_1 = \frac{p}{m\,(p^2+m^2)} .
\end{equation}
These are the photon wavefunctions in momentum space. They fall as $1/p^2$ and $1/(mp)$. If the
amplitude were just $\psi$, the cross section would fall as $p^{-4}$ (as in the inclusive case).

\section*{4. A factor $r^2$ is a derivative in $p$}

$J_\nu(pr)$ satisfies the Bessel equation. Written as a derivative in $p$:
\begin{equation}
  \Big(\frac{d^2}{dp^2} + \frac{1}{p}\frac{d}{dp} - \frac{\nu^2}{p^2}\Big) J_\nu(pr) = -\,r^2\, J_\nu(pr).
\end{equation}
The operator on the left is the 2D Laplacian $\nabla_p^2$ acting on a function with angular
dependence $e^{i\nu\phi}$. So
\begin{equation}
  \int dr\, r\, J_\nu(pr) K_\nu(mr)\; r^2 = -\Big(\frac{d^2}{dp^2} + \frac{1}{p}\frac{d}{dp} - \frac{\nu^2}{p^2}\Big)\psi_\nu(p).
\end{equation}

\textbf{$\nu=0$.} With $D\equiv p^2+m^2$:
\[
  \frac{d}{dp}\frac1D = -\frac{2p}{D^2},\qquad
  \frac1p\frac{d}{dp}\Big(p\cdot\frac{-2p}{D^2}\Big) = -\frac{4}{D^2}+\frac{8p^2}{D^3} = \frac{4(p^2-m^2)}{D^3}.
\]
\begin{equation}
  \boxed{\,I_0 = -\,Q_s^2\,\frac{p^2-m^2}{(p^2+m^2)^3}\,}\;\;\longrightarrow\;\; -\frac{Q_s^2}{p^4}.
  \label{eq:I0}
\end{equation}

\textbf{$\nu=1$.} With $f=p/D$:
\[
  f'' = -\frac{6p}{D^2}+\frac{8p^3}{D^3},\qquad \frac{f'}{p} = \frac{1}{pD}-\frac{2p}{D^2},\qquad \frac{f}{p^2}=\frac{1}{pD},
\]
\[
  f''+\frac{f'}{p}-\frac{f}{p^2} = -\frac{8p}{D^2}+\frac{8p^3}{D^3} = -\frac{8m^2p}{D^3}.
\]
\begin{equation}
  \boxed{\,I_1 = Q_s^2\,\frac{2\,m\,p}{(p^2+m^2)^3}\,}\;\;\longrightarrow\;\; \frac{2\,m\,Q_s^2}{p^5}.
  \label{eq:I1}
\end{equation}

Note the difference between the two. $\psi_0\sim 1/p^2$ and the Laplacian gives $1/p^4$, two more
powers, as expected. But $\psi_1\sim 1/(mp)$ and the Laplacian does \emph{not} give $1/p^3$: the
operator with $\nu=1$ acting on $1/p$ is exactly zero ($2-1-1=0$). Only the correction of order
$m^2/p^2$ survives, so $I_1\sim 1/p^5$. The reason: for $m=0$ the $\nu=1$ wavefunction is
$\vec p/p^2=\vec\nabla\ln p$, which is a harmonic function, $\nabla^2(\vec p/p^2)=0$ for $p\neq0$.

I checked Eqs.~\eqref{eq:I0} and \eqref{eq:I1} numerically (direct integration with $N=r^2/4$,
$m=1.5$, $p=3,6,12$~GeV): they agree to 10 digits.

\section*{5. The result}

\begin{equation}
  (z^2+\bar z^2)\,I_1^2 + I_0^2 \;=\; Q_s^4\;\frac{(p^2-m^2)^2 + 4\,(z^2+\bar z^2)\,m^2p^2}{(p^2+m^2)^6}
  \;\;\longrightarrow\;\; \frac{Q_s^4}{p^8}\qquad (p\gg m).
  \label{eq:g}
\end{equation}
With the prefactor of Eq.~\eqref{eq:code}:
\begin{equation}
  \frac{d\sigma}{d^2p\,dy}\;\sim\; z\,\frac{m^2\,Q_s^4}{p^8}.
\end{equation}
\begin{itemize}
\item The $p^{-8}$ is the $I_0$ term (the $K_0$, same-helicity part). The $I_1$ term is smaller by
$4(z^2+\bar z^2)\,m^2/p^2$.
\item The cross section is proportional to $m^2$: for massless quarks it vanishes at this order.
\item It is proportional to $Q_s^4$, not $Q_s^2$: two gluons in the amplitude, four in the cross section.
\item For $z=1/2$ the numerator is simply $p^4+m^4$.
\item $I_0$ changes sign at $p=m$. Near $p_c\simeq m_c$ the $I_0$ term has a zero and the $I_1$ term
fills it.
\end{itemize}

\section*{6. The physical reason}

Write the dipole amplitude with the transverse momentum $k$ given by the target,
\begin{equation}
  N(r) = \int d^2k\;\varphi(k)\,\big(1-e^{i\vec k\cdot\vec r}\big).
\end{equation}
The ``1'' is the two gluons attached to the same quark, the $e^{i\vec k\cdot\vec r}$ is one gluon on
the quark and one on the antiquark. In momentum space the amplitude is then
\begin{equation}
  A(\vec p) = \int d^2k\;\varphi(k)\,\big[\psi(\vec p)-\psi(\vec p-\vec k)\big].
  \label{eq:diff}
\end{equation}
The exchange has no colour and (after the integral over the impact parameter) no net momentum, so
the pair cannot tell the two attachments apart when $k\ll p$: they cancel. Expanding in $k/p$:
\[
  \psi(\vec p)-\psi(\vec p-\vec k) = \vec k\cdot\vec\nabla\psi - \tfrac12 k_ik_j\,\partial_i\partial_j\psi+\dots
\]
The first term vanishes in the angle average of $\vec k$. The second gives $-\tfrac14k^2\nabla^2\psi$. So
\begin{equation}
  A(\vec p)\simeq -\frac14\Big[\int d^2k\,k^2\varphi(k)\Big]\nabla^2\psi(\vec p) = -\frac{Q_s^2}{4}\,\nabla^2\psi(\vec p),
\end{equation}
which is exactly Section 4 (the same expansion of $1-e^{i\vec k\cdot\vec r}$ gives $N\simeq Q_s^2r^2/4$).

Compare:
\begin{center}
\begin{tabular}{lll}
\toprule
process & why & power \\
\midrule
inclusive $c\bar c$ & the target takes the recoil, no cancellation & $\sim Q_s^2/p^4$ \\
exclusive $c\bar c$ & cancellation, amplitude $\sim Q_s^2\nabla^2\psi$ & $\sim m^2Q_s^4/p^8$ \\
diffractive $c\bar c g$ & the extra gluon takes the recoil of the pair & hard factor $H\sim 1/p^4$ \\
\bottomrule
\end{tabular}
\end{center}
For the diffractive part, the hard factor in \texttt{diffractive\_kernel} is
$H = z\,[(z^2+\bar z^2)(p^4+m^4)+2m^2p^2]/(p^2+m^2)^4\to z(z^2+\bar z^2)/p^4$. This is why the
diffractive part is above the exclusive one at high $p_T$.

\textbf{When the expansion fails.} The expansion needs $\int d^2k\,k^2\varphi(k)$ to converge, that
is, no gluons with $k\sim p$. This is true for GBW (Gaussian $\varphi$). For a perturbative tail
$\varphi\sim 1/k^4$ the integral has a logarithm, $N\sim r^2\ln(1/r)$, and for BK
$N\sim r^{2\gamma}$ with $\gamma<1$. Then hard gluons with $k\sim p$ contribute directly in
Eq.~\eqref{eq:diff}. The effect is large for $I_1$, because its leading $r^2$ term was zero:
\begin{equation}
  N=\frac{Q_s^2}{4}\,r^{2\gamma}:\qquad I_1\simeq\frac{Q_s^2}{4m}\int dr\,r^{2\gamma}J_1(pr)
  = \frac{Q_s^2}{4m}\;\frac{4^{\gamma}\,\Gamma(1+\gamma)}{\Gamma(1-\gamma)}\;\frac{1}{p^{1+2\gamma}} ,
\end{equation}
which is zero for $\gamma=1$ (because $1/\Gamma(0)=0$) but not for $\gamma<1$. In the same way, for
$N=\tfrac{Q_s^2}{4}r^2\ln(1/r)$ one gets $I_1\simeq Q_s^2/(2mp^3)$ (checked numerically). Then
\begin{equation}
  m^2 I_1^2 \sim \frac{Q_s^4}{p^{2+4\gamma}}\quad(\text{no } m^2 \text{ left}),\qquad
  m^2 I_0^2\sim \frac{m^2Q_s^4}{p^{4+4\gamma}}\times\text{logs}.
\end{equation}
So at very large $p$ the $I_1$ term wins and the power is $-(2+4\gamma)$, which is $-6$ for
$\gamma\to1$. The pure $p^{-8}$ is the result for a dipole that is exactly $r^2$.

In our dipole (MV$^e$ at $x=0.01$) the logarithm is $\ln\!\big(1/(r\Lambda)+e_c\,e\big)$ with
$e_c\simeq 19$, so it is almost constant for $r\gtrsim 0.1$~GeV$^{-1}$: the initial condition is
close to $r^2$ in our $p_T$ range. From the tables in \texttt{data/Pb/mve/} I get
$d\ln N/d\ln r^2 = 0.94$--$0.96$ at $x=0.01$, and $0.83$--$0.90$ after two units of BK evolution
($x\simeq 1.4\times10^{-3}$). So the $I_1$ piece is not negligible when $x_P<0.01$ (table in Section 7).

\section*{7. In our range $p$ is not much larger than $m$}

Even for exact $r^2$, Eq.~\eqref{eq:g} reaches $p^{-8}$ slowly. For $z=1/2$ the local power is
\begin{equation}
  \frac{d\ln}{d\ln p}\,\frac{p^4+m^4}{(p^2+m^2)^6} = \frac{4p^4}{p^4+m^4}-\frac{12\,p^2}{p^2+m^2},
\end{equation}
which for $m=1.5$~GeV gives
\begin{center}
\begin{tabular}{lccccc}
\toprule
$p_c$ [GeV] & 3 & 5 & 8 & 12 & 20 \\
\midrule
local power & $-5.8$ & $-7.0$ & $-7.6$ & $-7.8$ & $-7.9$ \\
\bottomrule
\end{tabular}
\end{center}
It is always flatter than $-8$ and comes from above.

With the real dipole tables (\texttt{data/Pb/mve/glauber\_mve\_0}), at fixed $x$ and $z=1/2$, the local
power of $I_0^2+\tfrac12I_1^2$ is:
\begin{center}
\begin{tabular}{lccccc}
\toprule
$p_c$ range [GeV] & 3--4 & 4--6 & 6--8 & 8--12 & 12--16 \\
\midrule
$x=0.01$ ($Y=0$)              & $-6.2$ & $-7.0$ & $-7.4$ & $-7.5$ & $-7.4$ \\
$x=1.4\times10^{-3}$ ($Y=2$)  & $-5.2$ & $-6.0$ & $-6.6$ & $-6.8$ & $-6.8$ \\
$x=1.8\times10^{-4}$ ($Y=4$)  & $-4.7$ & $-5.4$ & $-6.0$ & $-6.3$ & $-6.4$ \\
\bottomrule
\end{tabular}
\end{center}
and the ratio $\tfrac12I_1^2/I_0^2$ at $p_c=8$~GeV is $0.16$, $0.67$ and $1.3$ for $Y=0,2,4$.
So the hard part alone is between $p^{-6}$ and $p^{-7.5}$ in our range, and it gets flatter with
evolution, as Section 6 says. (The proton table \texttt{data/proton/mve.dat} gives almost the same numbers.)

\section*{8. What the full calculation gives}

Local power $d\ln\sigma/d\ln p_{D^0}$ from the output files
(\texttt{output/0n0n/central\_values/BCFY/}, Pb, only the sample $b_d=0$):
\begin{center}
\begin{tabular}{lccccc}
\toprule
$p_{D^0}$ range [GeV] & 3--4 & 4--6 & 6--8 & 8--10 & 10--12 \\
\midrule
exclusive, $y=-2$ & $-7.6$ & $-8.2$ & $-8.3$ & $-8.5$ & $-8.4$ \\
exclusive, $y=0$  & $-7.5$ & $-8.6$ & $-9.3$ & $-9.6$ & $-9.7$ \\
exclusive, $y=+2$ & $-8.1$ & $-9.4$ & $-10.6$ & $-11.3$ & $-11.9$ \\
\midrule
diffractive, $y=-2$ & $-5.0$ & $-5.5$ & $-6.0$ & $-6.4$ & $-6.8$ \\
diffractive, $y=0$  & $-4.9$ & $-5.4$ & $-5.7$ & $-6.1$ & $-6.2$ \\
diffractive, $y=+2$ & $-5.8$ & $-6.6$ & $-7.4$ & $-7.8$ & $-8.3$ \\
\bottomrule
\end{tabular}
\end{center}
(The exclusive points above 3~GeV use 2000 VEGAS calls, so these numbers have some noise.)

The exclusive result is around $-8$ to $-9$, and about 3 powers steeper than the diffractive one, but
it is \textbf{steeper} than the hard part of Section 7. Three things in the full calculation make it steeper:
\begin{enumerate}
\item \textbf{Fragmentation.} $p_c=p_{D^0}/z_h>p_{D^0}$. A pure power law keeps its power after the
$z_h$ integral, but our hard part is not a pure power: the $D^0$ at $p_{D^0}$ sees the hard part at
a larger $p_c$, where it is steeper (table in Section 7).
\item \textbf{$x_P$ dependence.} $x_P=M_{q\bar q}^2/W^2 = m_T e^{-y}/(\sqrt{s}\,\bar z)$ grows with
$p_T$. Larger $x_P$ means less evolution, a smaller $Q_s^4$ and a steeper hard part (Section 7), until
$x_P$ reaches \texttt{XPO\_FROZEN}$=0.01$, where the dipole stops changing. Above
\texttt{XPO\_MAX}$=0.1$ the contribution is cut.
\item \textbf{Photon flux.} The photon needs $q^+\geq p^+=m_Te^{y}/\sqrt2$, which grows with $p_T$,
and the flux falls with the photon energy. This is the strongest at $y=+2$ (highest photon energy), where the
power goes to $-12$, and the weakest at $y=-2$, where the result stays close to $-8.4$.
\end{enumerate}
The $y$ dependence in the table agrees with points 2 and 3. I did not separate the size of the three
effects one by one.

\section*{Summary}

\begin{itemize}
\item Mechanism: colour-singlet exchange $\Rightarrow$ amplitude $=\psi(p)-\psi(p-k)\simeq-\tfrac14Q_s^2\nabla^2\psi$
$\Rightarrow$ $A\sim Q_s^2/p^4$ $\Rightarrow$ $\sigma\sim m^2Q_s^4/p^8$.
\item Exact for $N=Q_s^2r^2/4$: Eq.~\eqref{eq:g}.
\item $-8$ is the asymptotic power of the hard part for an $r^2$ dipole. In our range the hard part is
flatter ($-6$ to $-7.5$) and the flux, $x_P$ and fragmentation make the final result steeper
($-7.5$ to $-9.7$ at $y=0$).
\end{itemize}

\end{document}
